\subsection{Ecodistretto di Porto Marghera description}
The Ecodistretto di Porto Marghera was born in 2017, after the interruption of the former waste-to-energy plant. Its creation was driven by two different companies, Ecoprogetto Venezia S.r.l. and Eco-Ricicli Veritas, which converged into one in 2022, thus becoming Eco+Eco S.r.l. . This company forms part of the Veritas group, a municipally owned S.p.A. and one of the largest multiutility in Italy. The former Ecoprogetto Venezia S.r.l. was an integrated facility for processing unsorted waste; meanwhile, Eco-Ricicli Veritas operated as a platform for sorting waste from separate collection. Now, the two plants are still divided, but together they provide an environmentally-safe management of waste for approximately 890.000 people (Gruppo Veritas, n.d.). The Ecodistretto also welcomes third-party companies engaged in recycling materials at the early stages of the value chain.

\subsection{Selected Companies}
 In this work, only three companies within Ecodistretto di Porto Marghera are considered, with the aim of making the modeling process easier, albeit without providing a full-size analysis of the area. The companies modeled in this work are: 
 \begin{enumerate}[leftmargin=*]
    \item \textbf{Veritas}: it works in the whole Province of Venice and in the municipality of Mogliano Veneto (TV). In this work, Veritas includes only the former Eco-ricicli Veritas, responsible for waste collection and separation, and Metalrecycling Venice, which instead receives metal waste and recycles it. Thus, in this work its only output is \textcolor{sapienza}{recycled metal}. 
    \item \textbf{Ecopatè}: it is a glass recycling plant located in Musile di Piave. Recently, the plant was bought by Sibelco S.p.A which has completely substituted Ecopatè. Its only output is \textcolor{sapienza}{recycled glass}. 
    \item \textbf{Rilegno}: it is the Italian national consortium responsible for the collection and recycling of wooden packaging. It was established in 1997 following the D. Lgs No. 22/1997 (Decreto Ronchi). Its only output is \textcolor{sapienza}{recycled wood products}. 
 \end{enumerate}